\subsection{Charged and paraconsistent support}

\begin{theorem}[Genuine lax slack needs a non-thin carrier]\label{thm:pre-value-category}\lean{Infomaterial.Positing.Value.genuine_lax_cone_needs_non_thin_carrier}\leanok\uses{thm:value-category, thm:four-bilattice}
The value-layer generalization behind paraconsistent cones. \texttt{LaxSlackDiagnosis} shows that on
an antisymmetric (thin) carrier the KZ inequality collapses to strict idempotence, so a thin
\texttt{ValueCategory} carries \emph{no} genuine lax slack. A \texttt{PreValueCategory} drops
antisymmetry (a preorder with seed), with a lax closure satisfying $op(op\,v)\le op\,v$ instead of
idempotence. The structural theorem: if such a closure has real slack at any value
($op(op\,v)\le op\,v$ but $op(op\,v)\neq op\,v$), the carrier is non-thin --- there exist distinct,
mutually-$\le$ values. On an antisymmetric carrier the lax closure is idempotent, recovering the
thin case. This is added alongside, not replacing, \texttt{ValueCategory}.
\end{theorem}

\begin{theorem}[Paraconsistent causal reports]\label{thm:paraconsistent-cone}\lean{Infomaterial.Relational.Spacetime.Causal.ParaconsistentCone.paraconsistent_cone_reports}\leanok\uses{thm:cone-field-geometry, thm:pre-value-category}
Causal \emph{knowledge}, not physical spacetime, can be paraconsistent: a report assigns a
\texttt{Charge} to each ordered event pair, asserting reachability ($t$), its negation, both
(inconsistent $\top$), or neither (unknown $\bot$). Reports merge commutatively under the knowledge
join, and merging accumulates knowledge --- each cell climbs the knowledge order. A Boolean causal
relation embeds as a \emph{consistent} report (never $\top$) whose reachability verdict is exactly
the Boolean relation. Paraconsistency lives strictly in the report/observation layer --- the
substrate for contradictory light observations, agent traces, and conflicting causal
reconstructions; the non-thin value layer it requires is \cref{thm:pre-value-category}.
\end{theorem}

\begin{theorem}[Charged microsupport --- the cone valued in the FOUR]\label{thm:charged-microsupport}\lean{Infomaterial.Relational.Spacetime.Support.Charged.charged_microsupport}\leanok\uses{thm:paraconsistent-cone, thm:causal-microsupport}
The cross-faculty theorem the light-cone slate was built toward: the cone (\texttt{Lorentzian}) and
the FOUR (\texttt{Positing}) are \emph{one} structure. \cref{thm:causal-microsupport} propagated a
two-valued section; here the microsupport itself is lifted to the bilattice. A charged source at $a$
\textbf{emits} its charge into its causal future (the settled Boolean indicator $\mathrm{le}$) and
$\bot_k$ elsewhere, and two sources merge by the knowledge join. Then: (i) \textbf{spacelike /
acausal $\Rightarrow \bot$} --- a source emits the under-determined $\bot_k$ to every event outside
its cone, so spacelike-separated regions are knowledge-independent (told nothing); (ii)
\textbf{causal collision of contradiction $\Rightarrow \top$} --- when two influences carrying
\emph{contradictory} charges $t$ (affirm) and $f$ (deny) meet in a \textbf{common future event} (the
tip of a causal diamond), their join is the over-determined $\top_k$: the paraconsistent conflict
arises \emph{geometrically}, from causal convergence, and is \textbf{held} without explosion; (iii)
\textbf{agreement $\Rightarrow$ reinforcement} --- same-charge influences merge consistently, no
spurious conflict; (iv) merging \textbf{climbs $\leq_k$}. Spacelike incomparability is $\bot$, causal
order is $t$/$f$, held contradiction is $\top$: the light cone is the FOUR met as
causal admissibility. Axiom-free.
\end{theorem}

\begin{theorem}[Causal polarization --- the symplectic split as a proper colouring]\label{thm:causal-polarization}\lean{Infomaterial.Relational.Spacetime.Support.Polarization.causal_polarization}\leanok\uses{thm:causal-cochain-complex, thm:charged-microsupport, thm:colourings-pull-back, thm:four-colour-palette-is-four, thm:odd-cycle-obstructs-two-colouring}
The symplectic structure viewed through the colouring lens, with MLD as mediator --- the
polarization layer of \cref{thm:cotangent-symplectic-flow}. A symplectic structure needs a
\textbf{polarization} (a split into conjugate position/momentum directions), the antisymmetric sign
the bare directed-acyclic causal edges lack. The colouring supplies it: the polarization 1-form of a
colouring $c$ is its coboundary $\delta^0 c(a,b) = c(b) - c(a)$, and it is \textbf{nondegenerate}
(nonzero on every causal edge) \emph{exactly when} $c$ is a \textbf{proper colouring}
(\texttt{nondegenerate\_iff\_proper}) --- a degenerate causal direction is an edge whose endpoints
share a colour. At four colours the palette is the FOUR (\texttt{GraphColoring.fourColouring}), so the
conjugate split is valued in the Belnap bilattice, its two directions the two orders $\leq_k$,
$\leq_t$. The mediator: an \textbf{MLD} is a \textbf{factor map} (a graph homomorphism / sliding-block
code), and colourings \textbf{pull back along it} (\texttt{ProperColoring.comp}), so a factor map
carries a nondegenerate polarization (\texttt{mld\_preserves\_polarization}) --- MLD is the discrete
\textbf{canonical transformation}. And the \textbf{obstruction} to a global (Bool) polarization is
exactly \textbf{non-bipartiteness} --- an odd causal cycle: the chromatic $H^1$-class, the
MLD-invariant symplectic-curvature obstruction. This is the polarization layer and its obstruction;
together with the Liouville form it is the full discrete symplectic structure --- exact, not a
nondegenerate 2-class --- assembled in \cref{thm:cotangent-symplectic-flow}.
\end{theorem}
